%-------------------------
% Resume in Latex
% Author : Jake Gutierrez
% Based off of: https://github.com/sb2nov/resume
% License : MIT
%------------------------

\documentclass[letterpaper,11pt]{article}

\usepackage{latexsym}
\usepackage[empty]{fullpage}
\usepackage{titlesec}
\usepackage{marvosym}
\usepackage[usenames,dvipsnames]{color}
\usepackage{verbatim}
\usepackage{enumitem}
\usepackage[colorlinks=true, linkcolor=blue, urlcolor=blue, citecolor=blue]{hyperref}
\usepackage{fancyhdr}
\usepackage[english]{babel}
\usepackage{tabularx}
\input{glyphtounicode}
\usepackage{fontawesome5}
\usepackage{lmodern}

%----------FONT OPTIONS----------
% sans-serif
% \usepackage[sfdefault]{FiraSans}
% \usepackage[sfdefault]{roboto}
% \usepackage[sfdefault]{noto-sans}
% \usepackage[default]{sourcesanspro}

% serif
% \usepackage{CormorantGaramond}
% \usepackage{charter}


\pagestyle{fancy}
\fancyhf{} % clear all header and footer fields
\fancyfoot{}
\renewcommand{\headrulewidth}{0pt}
\renewcommand{\footrulewidth}{0pt}

% Adjust margins
\addtolength{\oddsidemargin}{-0.5in}
\addtolength{\evensidemargin}{-0.5in}
\addtolength{\textwidth}{1in}
\addtolength{\topmargin}{-.5in}
\addtolength{\textheight}{1.0in}

\urlstyle{same}

\raggedbottom
\raggedright
\setlength{\tabcolsep}{0in}

% Sections formatting
\titleformat{\section}{
  \vspace{-4pt}\scshape\raggedright\large
}{}{0em}{}[\color{black}\titlerule \vspace{-5pt}]

% Ensure that generate pdf is machine readable/ATS parsable
\pdfgentounicode=1

%-------------------------
% Custom commands
\newcommand{\resumeItem}[1]{
  \item\small{
    {#1 \vspace{-2pt}}
  }
}

\newcommand{\resumeSubheading}[4]{
  \vspace{-2pt}\item
    \begin{tabular*}{0.97\textwidth}[t]{l@{\extracolsep{\fill}}r}
      \textbf{#1} & #2 \\
      \textit{\small#3} & \textit{\small #4} \\
    \end{tabular*}\vspace{-7pt}
}

\newcommand{\resumeSubSubheading}[2]{
    \item
    \begin{tabular*}{0.97\textwidth}{l@{\extracolsep{\fill}}r}
      \textit{\small#1} & \textit{\small #2} \\
    \end{tabular*}\vspace{-7pt}
}

\newcommand{\resumeProjectHeading}[2]{
    \item
    \begin{tabular*}{0.97\textwidth}{l@{\extracolsep{\fill}}r}
      \small#1 & #2 \\
    \end{tabular*}\vspace{-7pt}
}

\newcommand{\resumeSubItem}[1]{\resumeItem{#1}\vspace{-4pt}}

\renewcommand\labelitemii{$\vcenter{\hbox{\tiny$\bullet$}}$}

\newcommand{\resumeSubHeadingListStart}{\begin{itemize}[leftmargin=0.15in, label={}]}
\newcommand{\resumeSubHeadingListEnd}{\end{itemize}}
\newcommand{\resumeItemListStart}{\begin{itemize}}
\newcommand{\resumeItemListEnd}{\end{itemize}\vspace{-5pt}}

%-------------------------------------------
%%%%%%  RESUME STARTS HERE  %%%%%%%%%%%%%%%%%%%%%%%%%%%%

\begin{document}
\vspace{-80pt}
%----------HEADING----------
\begin{center}
    \textbf{\Huge \scshape Ajay Kumar Patro} \\ \vspace{1pt}
    \small
    \faPhone\ +91 7894464335 \quad | \quad
    \faEnvelope\ \href{mailto:ajay789patro@outlook.com}{ajay789patro@outlook.com} \quad | \quad
    \faLinkedin\ \href{https://www.linkedin.com/in/ajay-patro-933659190/}{linkedin} \quad | \quad
    \faGithub\ \href{https://github.com/ajay098patro}{github.com/ajay098patro} \\ \vspace{2pt}
    GIS Backend \& Geospatial Data Engineer \quad | \quad PostGIS, Python, GeoServer, Keycloak
\end{center}
\vspace{-16pt}

 
%-----------PROGRAMMING SKILLS-----------
\section{Technical Skills}
 \begin{itemize}[leftmargin=0.15in, label={}]
    \small{\item{
     \textbf{Languages:}{ Python, SQL, JavaScript, Bash} \\
     \textbf{Frameworks:}{ Django, FastAPI} \\
     \textbf{Databases:}{ PostgreSQL/PostGIS, MySQL, SQLite} \\
     \textbf{Database Engineering:}{ Spatial SQL, Query Optimization, Indexing, Schema Design, Database Migrations} \\
     \textbf{GIS:}{ \\
        \hspace{1em}\textbf{Tools:}{ ArcGIS Pro, ArcMap, QGIS, PyQGIS, ArcPy, GDAL/OGR, Orfeo ToolBox, Geopandas, Google Earth Engine, PDAL}\\
        \hspace{1em}\textbf{Services:}{ GeoServer, Martin, pg\_tileserv, TiTiler}
     } \\
     \textbf{Infrastructure \& Tools:}{ Docker, Docker Compose, Linux, IIS, Apache, Tomcat, Nginx, Git, GitHub} \\
     \textbf{Backup \& Auth:}{ pgBackRest (PITR), Cloudflare R2, Keycloak/OIDC}
    }}
 \end{itemize}

 \vspace{-19pt}


%-----------EXPERIENCE-----------      
\section{Experience}
  \resumeSubHeadingListStart
\resumeSubheading
{GIS Analyst}{Mar 2026 -- Present}
{Monarch Surveyors \& Engineering Consultants Ltd}{Nashik, Maharashtra, India}

\resumeItemListStart
\resumeItem{Deployed a 2-node \textbf{GeoServer} cluster for a state government agency behind \textbf{IIS} (reverse proxy) with \textbf{Apache Tomcat}, \textbf{Keycloak}/OIDC and HTTPS/SSL, using an \textbf{SMB} share over the existing \textbf{SAN} as a common data directory.}
\resumeItem{Wrote a \textbf{batch script} to deploy primary and secondary nodes, remove existing installations, and promote either node to primary, enabling manual failover in 10--20 minutes, within the client's 1-hour downtime tolerance.}
\resumeItem{Designed and set up PostgreSQL/PostGIS backup and recovery from scratch using \textbf{pgBackRest} with \textbf{Cloudflare R2}, with weekly full backups and \textbf{point-in-time recovery (PITR)} for a $\sim$9\,GB database; adopted by the team.}
\resumeItem{Managed \textbf{PostgreSQL/PostGIS} and \textbf{GeoServer} for a multi-tenant LRMS platform supporting government land acquisition and infrastructure projects, including spatial schema design and ingestion, transformation and validation of survey, drawing and \texttt{.shp} datasets.}
\resumeItem{Converted drone imagery from \texttt{.ecw} to Cloud Optimized GeoTIFFs (COGs) with internal overviews and quantized alpha bands, and published raster and vector layers through \textbf{GeoServer} (WMS/WFS/ImageMosaic) for dashboard web maps.}
\resumeItem{Fixed geometry, SRID and topology problems in \textbf{PostGIS} and debugged spatial SQL, raster rendering and data-sync issues across \textbf{PostGIS} and \textbf{GeoServer}.}
\resumeItemListEnd
      
      
      \resumeSubheading
  {GIS Specialist (Python Data Ops Automation)}{Apr 2024 -- May 2025}
  {Lepton Software Export \& Research Pvt. Ltd.}{Pune, Maharashtra, India}

\resumeItemListStart
  \resumeItem{\textbf{Project:} GOFNMS Tele Infra (telecom network data processing)}
  \resumeItem{Populated and maintained fibre-network records in the project database from stakeholder data: fibre direction, inside-plant (ISP) data, and subscriptions mapping each fibre to its \textbf{OTU/IPS} port.}
  \resumeItem{Reconciled live \textbf{OTU/IPS} fibre data against the database, reviewed fibre-length mismatches flagged by \textbf{Enghouse NetDesigner} (ArcGIS-based), and reported them to the team responsible for corrections, supporting fault analysis.}
  \resumeItem{Created inside-plant (ISP) records from field network engineers' reports to support telecom network inventory and planning.}
  \resumeItem{Built Python (\textbf{ArcPy}) automation for daily spatial data processing, validation, and reporting.}
\resumeItemListEnd

      \resumeSubheading
  {Scientific Intern (C-GARD, NIRDPR)}{Mar 2023 -- Jul 2023}
  {NIRD \& PR NERC}{Guwahati, Assam}
\resumeItemListStart
\resumeItem{Maintained and validated geospatial datasets for infrastructure development planning, including ground assessments aligned with \textbf{RADFI 2017} guidelines.}
\resumeItem{Delivered technical tutorials on the \textbf{Yuktadhara--Bhuvan Portal} to field officers across Assam and Meghalaya.}
\resumeItem{Acted as a secondary-level reviewer for research papers in the ISBN-published C-GARD book, \textit{\textbf{Geoinformatics Application in Land, Water, and Infrastructure Development}}.}
\resumeItemListEnd

\vspace{-10pt}
%-----------PROJECTS----------- 
\section{Projects} 
\vspace{-10pt}
\resumeProjectHeading
{\textbf{Secure Vector-Tile Serving Platform: Martin + APISIX + Keycloak (In Progress)}}{}
\resumeItemListStart
\resumeItem{\textbf{Tech Stack:}\textit{ Martin, PostgreSQL/PostGIS, Docker Compose, Apache APISIX, Keycloak/OIDC}}
\resumeItem{Building a containerized service that serves vector tiles from \textbf{PostGIS} through \textbf{Martin}, with \textbf{Apache APISIX} enforcing \textbf{Keycloak}-based role and resource authorization before requests reach the tile server.}
\resumeItem{Designing cached authorization decisions to keep per-request validation fast, and working out how cached tiles and decisions should stay correct when data or roles change.}
\resumeItemListEnd

\resumeProjectHeading 
{\textbf{End-to-End AI Burn Scar Detection Pipeline (Ongoing)}}{} 
\resumeItemListStart 
\resumeItem{\textbf{Tech Stack:}\textit{ PyTorch, Apache Airflow, STAC Catalog, FastAPI, PostGIS, GeoServer, Rasterio, GeoPandas}} 
\resumeItem{Developed an automated pipeline to ingest Sentinel-2 imagery via \href{https://dataspace.copernicus.eu/}{CDSE's} \textbf{STAC API} to detect wildfire burn scars using the \href{https://huggingface.co/ibm-nasa-geospatial/Prithvi-EO-2.0-300M-TL}{HuggingFace Model: ibm-nasa-geospatial Prithvi EO-2.0} Model.} 
\resumeItem{Orchestrated end-to-end \textbf{ETL} workflows with Apache Airflow for preprocessing, inference, and spatial data loading into \textbf{PostGIS}.} 
\resumeItem{Generated STAC-compliant metadata and published GIS-ready spatial outputs via \textbf{GeoServer}, exposing them through \textbf{FastAPI} as REST API endpoints for downstream analytics and visualization.} 
\resumeItemListEnd 

\resumeProjectHeading 
{\textbf{Spectral--Index Feature Engineering / OBIA Pipeline}}{} 
\resumeItemListStart 
\resumeItem{Built an end-to-end \textbf{Object-Based Image Analysis (OBIA)} pipeline for satellite imagery, reducing land-cover mapping time by \textbf{$\sim$92\%} by deriving spectral-index--based object features for automated training label generation. Implemented using \textbf{QGIS Model Builder} and \textbf{Orfeo ToolBox}.}
\resumeItemListEnd

\resumeProjectHeading
{\textbf{LiDAR Point-Cloud Web Viewer (Potree, Range-Aware HTTP Streaming)}}
 {\faGithub\ \href{https://github.com/ajay098patro/lidar-potree-range-streaming}{Source Code}}
\resumeItemListStart
\resumeItem{Built a web-based LiDAR point-cloud viewer using \href{https://github.com/potree/potree}{Potree} for large-scale multi-million point visualization. Engineered a range-aware static HTTP streaming server to replace Django static serving, enabling level-of-detail (LOD) octree loading, eliminating repeated full-file downloads, and optimizing browser memory and rendering performance. Implemented using Django, Potree and Flask.}
\resumeItemListEnd

\resumeProjectHeading 
{\textbf{Football Field Booking Management System (CS50 SQL Capstone)}}
 {\faGithub\ \href{https://github.com/ajay098patro/Football-Field-Booking-Management-System-CS50-SQL-Final-Project}{Source Code}}
\resumeItemListStart 
\resumeItem{Designed a normalized, ACID-compliant \textbf{MySQL} schema (10+ tables) with constraints, triggers and indexing; tuned queries with EXPLAIN.} 
\resumeItemListEnd
\resumeSubHeadingListEnd
\vspace{-8pt}
%-----------EDUCATION-----------
\section{Education}
\small % keeps font readable and consistent
\setlength{\itemsep}{1pt} % minimal spacing between entries
\setlength{\parskip}{0pt}
\resumeSubHeadingListStart

\resumeSubheading
{Central University of Karnataka}{ Nov 2021 -- Aug 2023}{\hspace{1em}M.Sc. in Applied Geography \& Geoinformatics}{Kalaburagi, India}
\vspace{0.1pt}
\par
{\scriptsize \textbf{Coursework:} \textbf{GeoAI}, \textbf{LULC} classification (supervised and unsupervised), \textbf{hydrological} and basin morphometric analysis (hypsometric integral, \textbf{HRI/HRS}), soil erosion modeling (\textbf{USLE}, \textbf{MUSLE}), \textbf{geostatistics} (hotspot analysis, spatial autocorrelation), geocoding, reverse geocoding, and \textbf{multi-criteria decision analysis} (MCDA) for site suitability.}

\resumeSubHeadingListEnd

\vspace{-13pt}
%-----------CERTIFICATIONS-----------
\section{Certifications}
\begin{itemize}[leftmargin=0.03in, label={}]
\footnotesize
\setlength\itemsep{2pt} % compact spacing
\setlength\parskip{0pt} 

\item {\textbf{CS50 Python Certificate} \hfill 
      \faAward\ \href{https://cs50.harvard.edu/certificates/7ba41367-a931-47cb-ac47-ccba96f4048d}{View Certificate} \\
      \hspace{1em}{Harvard's OCW CS50, Cambridge, Massachusetts} \hfill 2023 -- 2024}

\item {\textbf{CS50 SQL Certificate}} \hfill 
      \faAward\ \href{https://cs50.harvard.edu/certificates/0f4666c7-f10d-430c-92e9-933fb996e84f}{View Certificate} \\
      \hspace{1em}{Harvard's OCW CS50, Cambridge, Massachusetts} \hfill 2025

\item {\textbf{ArcGIS Enterprise MOOC} \hfill \\
      \hspace{1em}{Esri} \hfill 2026 -- Present}

\end{itemize}

\end{document}
